\section{Заключение}

В работе продемонстрирована состоятельность методологии сквозного бескомпромиссного HW/SW-проектирования систем цифровой обработки навигационных сигналов с использованием открытых инструментов[cite: 13, 14]. Разработанный тракт параллельного поиска обеспечивает надежное обнаружение сигналов GPS L1 при низких отношениях сигнал/шум (до $-20\text{ дБ}$) и существенной динамике частоты[cite: 10, 11, 15]. 

Компиляция синтезируемого SystemVerilog через Verilator обеспечивает прирост скорости имитации на 2--3 порядка по сравнению с промышленными логическими симуляторами[cite: 14], формируя строгую базу для создания открытых RISC-V систем на кристалле космического базирования.
